\section{Formula dictionary}
\label{gq:app:formula-dictionary}

This appendix collects the principal transforms in one place.  It is meant
as a reference sheet rather than a substitute for the proofs.

\subsection{General geometric comb}

For $x_j=cq^j$, $0<q<1$, define
\[
 N_k(x)=\prod_{r=0}^{k-1}(x-cq^r),
 \qquad
 \omega_{n+1}(x)=N_{n+1}(x).
\]
Then
\begin{longtable}{@{}p{0.28\textwidth}p{0.64\textwidth}@{}}
\caption{Core geometric-comb formulas.}\label{gq:tab:formula-dictionary-general}\\
\toprule
Object & Formula\\
\midrule
\endfirsthead
\toprule
Object & Formula\\
\midrule
\endhead
Newton/$q$-Pochhammer basis
& $N_k(x)=(-c)^kq^{\binom k2}(x/c;q^{-1})_k$\\[4pt]
Divided difference
& $[f;c,\ldots,cq^k]=D_q^kf(c)/[k]_q!$\\[4pt]
Explicit data transform
& $c^{-k}\sum_{j=0}^k
 (-1)^jq^{-jk+\binom{j+1}{2}}f(cq^j)/
 ((q;q)_j(q;q)_{k-j})$\\[6pt]
Monomial to Newton
& $x^m=\sum_{k=0}^{m}c^{m-k}\GaussianBinomial{m}{k}{q}N_k(x)$\\[4pt]
Newton to monomial
& $N_k(x)=\sum_{m=0}^{k}(-1)^{k-m}c^{k-m}
 q^{\binom{k-m}{2}}\GaussianBinomial{k}{m}{q}x^m$\\[6pt]
Barycentric weight
& $\lambda_{j,n}=c^{-n}(-1)^jq^{\binom{j+1}{2}-nj}/
 ((q;q)_j(q;q)_{n-j})$\\[6pt]
Accumulation weight
& $\mu_{j,n}=(-1)^{n-j}q^{\binom{n-j+1}{2}}/
 ((q;q)_j(q;q)_{n-j})$\\[6pt]
Endpoint Lebesgue constant
& $\Leb_n(0)=(-q;q)_n/(q;q)_n$\\[4pt]
Global Lebesgue asymptotic
& $\Leb_n^*\sim2(-q;q)_\infty q^{-\binom n2}/(\EulerE n)$\\[4pt]
Monomial residual
& $x^{n+1+r}-\Interp_nx^{n+1+r}
 =N_{n+1}(x)h_r(x,c,cq,\ldots,cq^n)$\\[6pt]
Geometric specialization
& $h_r(c,cq,\ldots,cq^n)=c^r\GaussianBinomial{n+r}{n}{q}$\\[4pt]
Jackson moment
& $\int_0^cN_k(x)\,\Differential_qx=(-1)^kc^{k+1}
 q^{\binom{k+1}{2}}/[k+1]_q$\\
\bottomrule
\end{longtable}

\subsection{Half-base Fabius formulas}

For $q=1/2$, $d_n=\Expectation X^n$, and $a_n=d_n/n!$,
\begin{longtable}{@{}p{0.28\textwidth}p{0.64\textwidth}@{}}
\caption{Fabius and Rvachev specialization.}\label{gq:tab:formula-dictionary-fabius}\\
\toprule
Object & Formula\\
\midrule
\endfirsthead
\toprule
Object & Formula\\
\midrule
\endhead
Moment recurrence
& $(2^n-1)a_n=\sum_{j<n}a_j/(n-j+1)!$\\[5pt]
Rvachev half-moment
& $d_n=n\int_0^1t^{n-1}\RvachevUp(t)\,\Differential t$\\[5pt]
Inverse-dyadic value
& $F(2^{-n})=2^{-\binom n2}a_n$\\[5pt]
Newton coefficient
& $\FabiusGeometricNewtonCoefficient{k}{q}=(q;q)_k^{-1}\sum_{j=0}^k(-1)^j
 \GaussianBinomial{k}{j}{q^{-1}}a_j$\\[6pt]
Finite interpolant
& $\Interp_nF(x)=\omega_{n+1}(x)(q;q)_n^{-1}
 \sum_{j=0}^n(-1)^j\GaussianBinomial{n}{j}{q^{-1}}a_j/(x-q^j)$\\[7pt]
Boundary residue generating function
& $\sum_{n\ge0}R_nz^n=(qz;q)_\infty
 \sum_{j\ge0}q^{\binom j2}a_jz^j/(q;q)_j$\\[7pt]
Boundary bound
& $|R_n|\le\EulerE(q;q)_\infty^{-2}q^{\Floor{n^2/4}}$\\[5pt]
Fixed-jet bound
& $|(\Interp_nF)^{(m)}(0)|\le
 \EulerE m!\binom nm(q;q)_\infty^{-2}
 q^{\Floor{n^2/4}-mn+\binom m2}$\\[7pt]
Rvachev relation
& $\RvachevUp(t)=F(1-|t|)$ and $\Interp_n\RvachevUp=1-\Interp_nF$
 on the center-directed positive comb\\
\bottomrule
\end{longtable}

\subsection{Atomic synthesis formulas}

For a degree bound $n$, $h=2^{-n}$ and
$K_n=\{-2^{n+1}+1,\ldots,2^{n+1}-1\}$,
\begin{align}
 \App p&=\MGF_u(D)^{-1}p
 =\sum_{r\ge0}\frac{\gamma_{2r}}{(2r)!}p^{(2r)},
 \label{gq:eq:dictionary-App}\\
 p(x)&=h\sum_{k\in K_n}(\App p)(kh)u(x-kh),
 \qquad \deg p\le n,
 \label{gq:eq:dictionary-synthesis}\\
 \ell_{j,n}(x)&=h\sum_{k\in K_n}
 (\App\ell_{j,n})(kh)u(x-kh),
 \label{gq:eq:dictionary-cardinal-synthesis}\\
 (hB)D&=I_{n+1}.
 \label{gq:eq:dictionary-right-inverse}
\end{align}

\section{Selected exact dyadic data}
\label{gq:app:exact-data}

\subsection{Moments and inverse-dyadic values}

\begin{table}[htbp]
\centering
\caption{Regularized moments $a_n=d_n/n!$ and the corresponding values
$F(2^{-n})=2^{-\binom n2}a_n$.}
\label{gq:tab:exact-fabius-data}
\small
\begin{tabular}{@{}rll@{}}
\toprule
$n$ & $a_n$ & $F(2^{-n})$\\
\midrule
0 & $1$ & $1$\\
1 & $1/2$ & $1/2$\\
2 & $5/36$ & $5/72$\\
3 & $1/36$ & $1/288$\\
4 & $143/32400$ & $143/2073600$\\
5 & $19/32400$ & $19/33177600$\\
6 & $1153/17146080$ & $1153/561842749440$\\
7 & $583/85730400$ & $583/179789679820800$\\
8 & $1616353/2623350240000$
  & $1616353/704200217922109440000$\\
\bottomrule
\end{tabular}
\end{table}

\subsection{First geometric Newton coefficients}

For
\[
 \Interp_nF(x)=\sum_{k=0}^n\FabiusGeometricNewtonCoefficient{k}{1/2}
 \prod_{r=0}^{k-1}(x-2^{-r}),
\]
the exact coefficients begin
\begin{longtable}{@{}rl@{}}
\caption{First exact half-base Newton coefficients.}
\label{gq:tab:exact-newton-coefficients}\\
\toprule
$k$ & $\FabiusGeometricNewtonCoefficient{k}{1/2}$\\
\midrule
\endfirsthead
\toprule
$k$ & $\FabiusGeometricNewtonCoefficient{k}{1/2}$\\
\midrule
\endhead
0 & $1$\\
1 & $1$\\
2 & $-26/27$\\
3 & $-128/27$\\
4 & $-4253248/637875$\\
5 & $189673472/19774125$\\
6 & $134859284873216/1648153544625$\\
7 & $21867230161534976/209315500167375$\\
8 & $-349286360632157831954432/204161105975753390625$\\
\bottomrule
\end{longtable}

Large numerator and denominator sizes are not numerical artifacts; they are
exact consequences of composing the Fabius moment recurrence with the
reciprocal Gaussian transform.


\section{Additional derivations}

\subsection{Partial fractions for the $q$-exponential perpetuity}

For completeness, consider
\[
 L_q(s)=\prod_{j=0}^{\infty}(1+sq^j)^{-1}.
\]
At the simple pole $s=-q^{-k}$, the coefficient of
$(1+sq^k)^{-1}$ is
\begin{align*}
 c_k
 &=\left[\prod_{j\ne k}(1-q^{j-k})\right]^{-1}\\
 &=\frac{(-1)^kq^{k(k+1)/2}}{(q;q)_k(q;q)_\infty}.
\end{align*}
Indeed, the factors with $j<k$ give
$(-1)^kq^{-k(k+1)/2}(q;q)_k$, and the factors with $j>k$ give
$(q;q)_\infty$.  Since
\[
 \mathcal L^{-1}\left[\frac1{1+sq^k}\right](y)
 =q^{-k}e^{-y/q^k},
\]
we recover \eqref{g1:eq:Z-density}.

\subsection{Equivalence of the power formulas}

For an integer $d=n+1+r$, \eqref{g1:eq:pure-power} and
\eqref{g1:eq:all-moments} agree.  Indeed,
\begin{align*}
 q^{dn}\frac{(q^{1-d};q)_n}{(q;q)_n}
 &=q^{(n+1+r)n}
 \frac{\prod_{j=0}^{n-1}(1-q^{-n-r+j})}{(q;q)_n}\\
 &=(-1)^nq^{\binom{n+1}{2}}
 \frac{(q;q)_{n+r}}{(q;q)_n(q;q)_r}\\
 &=(-1)^nq^{\binom{n+1}{2}}\GaussianBinomial{n+r}{n}{q}.
\end{align*}
This calculation is often useful when moving between Mellin exponents and
integer Taylor moments.

\subsection{A generating function for the smooth remainder on exponentials}

Take $f(x)=e^{tx}$ in \eqref{g1:eq:HG-remainder}.  Using
\eqref{g1:eq:Y-moments},
\begin{align}
 (\Gc_{n,q}e^{t\cdot})(A)-1
 &=\frac{(-1)^nA^{n+1}q^{\binom{n+1}{2}}t^{n+1}}{(n+1)!}
 \Expectation e^{tY_{n,A,q}}\notag\\
 &=(-1)^nA^{n+1}q^{\binom{n+1}{2}}t^{n+1}
 \sum_{r=0}^{\infty}
 \frac{A^r\GaussianBinomial{n+r}{n}{q}}{(n+r+1)!}t^r.
 \label{g1:eq:exp-remainder}
\end{align}
The same series follows from \eqref{g1:eq:exact-analytic-tail}.  The equality is a
useful consistency check between the analytic-tail and B-spline pictures.


\section{Selected finite identities}
\label{g2:app:finite-identities}

For quick reference, the principal finite formulas are collected here.

\begin{align*}
 \Omega_{n+1}(x)
 &=x^{n+1}(a/x;q)_{n+1},\\[0.3em]
 \Phi_k(x)
 &=(-a)^kq^{\binom k2}(x/a;q^{-1})_k,\\[0.3em]
 \Phi_k(aq^j)
 &=(-a)^kq^{\binom k2}(q;q)_k\GaussianBinomial jkq,\\[0.3em]
 \beta_{n,j}
 &=\frac{(-1)^jq^{-jn+\binom{j+1}{2}}}
 {a^n(q;q)_j(q;q)_{n-j}},\\[0.3em]
 \lambda_{n,j}^{(q)}
 &=\frac{(-1)^{n-j}q^{\binom{n-j+1}{2}}}
 {(q;q)_j(q;q)_{n-j}},\\[0.3em]
 \sum_{j=0}^n\lambda_{n,j}^{(q)}z^j
 &=\frac{\prod_{r=1}^n(z-q^r)}{(q;q)_n},\\[0.3em]
 f[a,aq,\ldots,aq^k]
 &=\frac{D_q^kf(a)}{[k]_q!},\\[0.3em]
 \sum_{j=0}^n|\lambda_{n,j}^{(q)}|
 &=\frac{(-q;q)_n}{(q;q)_n},\\[0.3em]
 M_{n,q}(s)
 &=q^{ns}\frac{(q^{1-s};q)_n}{(q;q)_n},\\[0.3em]
 F_q(q^n)
 &=\frac{q^{\binom{n+1}{2}}}{(1-q)^nn!}\Expectation X_q^n
 \quad(0<q\le1/2),\\[0.3em]
 F(2^{-n})
 &=2^{-\binom n2}\frac{\mu_n}{n!},\\[0.3em]
 A_n
 &=\frac1{(1/2;1/2)_n}
 \sum_{j=0}^n(-1)^j\GaussianBinomial nj2\frac{\mu_j}{j!}.
\end{align*}

\section{Pseudocode for exact dyadic computation}
\label{g2:app:pseudocode}

The following high-level algorithm uses rational arithmetic throughout.

\begin{lstlisting}[style=pythonstyle]
mu[0] = 1
M[0]  = 1
for n = 1..N:
    mu[n] = (sum_{k=0}^{n-1} binom(n,k)*mu[k]/(n-k+1)) / (2^n-1)
    M[n]  = M[n-1]*(2^n-1)

for n = 0..N:
    Fdyadic[n] = 2^(-n*(n-1)/2) * mu[n] / n!
    A[n] = 2^(n*(n+1)/2) *
           sum_{j=0}^n (-1)^j * mu[j] / (j!*M[j]*M[n-j])
    E[n] = sum_{j=0}^n (-1)^(n-j) * 2^j * mu[j]
           / (j!*M[j]*M[n-j])
\end{lstlisting}

The attached Python implementation additionally verifies all transforms by
independent formulas, evaluates Newton partial sums at rational targets, and
computes the high-precision Fourier-product experiment.


\section{Gamma-product calculations}
\label{app:gamma}

\subsection{Half-cell central cardinal (mode A)}
In the scaled variable $y=Mx$ with interior nodes $1,\dots,M-1$ and
$j=M/2$, at $y=\tfrac12$:
$\prod_{k=1}^{M-1}\AbsoluteValue{k-\tfrac12}=\Gamma(M-\tfrac12)/\Gamma(\tfrac12)$;
deleting the $j$ factor divides by $\AbsoluteValue{j-\tfrac12}=\tfrac M2-\tfrac12$,
and the denominator of the cardinal is
$(j-1)!\,(M-1-j)!=\Gamma(M/2)^{2}$, giving $A_M$ in
\eqref{eq:A-mode}.  The weight ratio is
$\bigl[\tfrac1{2M}(1-\tfrac1{2M})\bigr]/\tfrac14=(2M-1)/M^{2}=a_M$.
Stirling gives $A_M\sim2^{M}/(\pi\sqrt2\,M)$.

\subsection{Near-central boundary-node cardinal (mode B)}
For $j=1$ at $y=(M+1)/2$ with $M=2K$:
$\AbsoluteValue{\prod_{k=2}^{M-1}(y-k)}
=\bigl[\Gamma(K-\tfrac12)/\Gamma(\tfrac12)\bigr]^{2}$
(the products on the two sides of the half-integer are equal), and
the denominator is $(M-2)!=\Gamma(M-1)$, giving $B_M$.  The weight
ratio is
$\bigl[(\tfrac12+\tfrac1{2M})(\tfrac12-\tfrac1{2M})\bigr]
/\bigl[\tfrac1M(1-\tfrac1M)\bigr]=(M+1)/4=b_M$.  The duplication
formula turns $B_M$ into
$\dfrac{2^{2-M}}{\sqrt\pi}\,
\dfrac{\Gamma((M-1)/2)}{\Gamma(M/2)}
\sim2^{2-M}\sqrt{2/(\pi M)}$; equivalently
$B_M=\binom{M-2}{M/2-1}\big/16^{\,M/2-1}$.

\subsection{The one-third mode}
At $y=M/3$ (dyadic $M$ is never divisible by $3$):
$\prod_{k=1}^{M-1}(M/3-k)=\Gamma(M/3)/\Gamma(1-2M/3)$, and reflection
gives $1/\AbsoluteValue{\Gamma(1-2M/3)}
=\Gamma(2M/3)\AbsoluteValue{\sin(2\pi M/3)}/\pi
=\tfrac{\sqrt3}{2\pi}\Gamma(2M/3)$.  Removing the central factor
$\AbsoluteValue{M/3-M/2}=M/6$ and dividing by $\Gamma(M/2)^{2}$ gives
\eqref{eq:one-third-mode}; the exponential part of the gamma ratio is
$\exp\{M[\tfrac13\log\tfrac13+\tfrac23\log\tfrac23-\log\tfrac12]\}
=\EulerE^{M(\log2-\Hyp(1/3))}$.

\subsection{Full-comb central cardinal}
For the full node set $0,\dots,M$ at $y=\tfrac12$, $j=M/2$:
$\prod_{k=0}^{M}\AbsoluteValue{\tfrac12-k}
=\tfrac12\,\Gamma(M+\tfrac12)/\Gamma(\tfrac12)$; deleting the
central factor divides by $\AbsoluteValue{\tfrac12-\tfrac M2}=\tfrac{M-1}2$,
and the denominator is
$(M/2)!\,(M/2)!=\Gamma(\tfrac M2+1)^{2}$, giving
\eqref{eq:full-central-cardinal}; Stirling's formula then yields
$\AbsoluteValue{\ell_{M/2}(\tfrac1{2M})}\sim\sqrt2\,2^{M}/(\pi M^{2})$, whose
square is \eqref{eq:hermite-square}.

\section{Mode asymptotics}
\label{app:mode-asymptotics}

Stirling's expansion in logarithmic form,
\begin{equation}
 \log\Gamma(z)
 =\Bigl(z-\tfrac12\Bigr)\log z-z+\tfrac12\log(2\pi)
 +\sum_{r=1}^{R-1}\frac{B_{2r}}{2r(2r-1)\,z^{2r-1}}
 +\BigO\bigl(z^{-2R+1}\bigr),
 \label{eq:stirling}
\end{equation}
applied to \eqref{eq:A-mode}--\eqref{eq:B-mode} yields
\begin{align}
 \log A_M&=M\log2-\log M-\log(\pi\sqrt2)+\BigO(M^{-1}),\\
 \log B_M&=-(M-2)\log2+\tfrac12\log\frac{2}{\pi M}+\BigO(M^{-1}),\\
 \log a_M&=-\log(M/2)+\log\bigl(1-\tfrac1{2M}\bigr),
 \qquad
 \log b_M=\log(M/4)+\log\bigl(1+\tfrac1M\bigr),
\end{align}
which prove \cref{prop:balance-asymptotic}.  Retaining further
Bernoulli terms and reverting with complete Bell polynomials yields
the all-orders program \eqref{eq:balance-all-orders}.

\section{Row reciprocity}
\label{app:row-symmetry}

With $A_m(z)=\sum_{k=0}^{M-1}F(k/M)z^{k}$ and
$F((M-k)/M)=1-F(k/M)$,
\begin{align}
 z^{M}A_m(1/z)
 &=\sum_{k=1}^{M}F\bigl((M-k)/M\bigr)z^{k}
 =\sum_{k=1}^{M}\bigl(1-F(k/M)\bigr)z^{k}
 =\frac{z(1-z^{M})}{1-z}-A_m(z)-z^{M},
\end{align}
which is \eqref{eq:AN-reciprocal}.  This exact reciprocity forces the
palindromic structure of the low-level row numerators
\eqref{eq:row-N1}--\eqref{eq:row-N3}.

\section{Core algorithms}
\label{app:algorithms}

\begin{lstlisting}[style=pythonstyle,caption={Streaming the exact
dyadic Fabius row (\cref{alg:stream}); \texttt{d} holds the exact
Fabius-law moments of \eqref{eq:d-recurrence}.}]
def fabius_row_exact(m):
    M = 2**m
    d = fabius_law_moments(m)          # Fractions d_0..d_m
    S = [0]*(m+1)                      # signed prefix power sums
    scale = Fraction(1, 2**(m*(m-1)//2) * factorial(m))
    row = []
    for a in range(M+1):
        row.append(scale * sum(comb(m,j)*d[j]*S[m-j]
                               for j in range(m+1)))
        if a == M: break
        eps = -1 if bin(a).count('1') % 2 else 1
        S = [S[0]+eps] + [sum(comb(k,q)*S[q] for q in range(k+1))
                          for k in range(1, m+1)]
    return row
\end{lstlisting}

\begin{lstlisting}[style=pythonstyle,caption={Endpoint-flat
interpolation in barycentric form (\cref{thm:factorization}).}]
S = lambda r,x: (factorial(2*r+1)//(factorial(r)**2)
                 * sum((-1)**k*comb(r,k)*x**(r+k+1)/(r+k+1)
                       for k in range(r+1)))
W = lambda r,x: (x*(1-x))**(r+1)
g = [(F(x_j)-S(r,x_j))/W(r,x_j) for x_j in interior_nodes]
w = [(-1)**(j-1)*comb(M-2,j-1) for j in range(1,M)]
H = S(r,x) + W(r,x)*barycentric(x, interior_nodes, g, w)
\end{lstlisting}

The absorbed packages add the sparse-difference row generation
($\BigO(mM)$), confluent divided differences for full-grid Hermite,
weighted-Lebesgue scans, and all figure/table generation; the canonical entry points, dependencies, and replay commands are
recorded in \nolinkurl{assets/README.md} and \nolinkurl{assets/VALIDATION.md}.


\section{Status ledger}
\label{gq:app:status-ledger}

\small
\begin{longtable}{@{}>{\raggedright\arraybackslash}p{0.27\textwidth}
>{\raggedright\arraybackslash}p{0.13\textwidth}
>{\raggedright\arraybackslash}p{0.335\textwidth}
>{\raggedright\arraybackslash}p{0.18\textwidth}@{}}
\caption{Logical and formal status of the principal claims.  ``Paper proof''
means a manuscript proof in this report, not a Lean declaration.}
\label{gq:tab:status-ledger}\\
\toprule
Claim & Manuscript & Lean boundary & Location\\
\midrule
\endfirsthead
\toprule
Claim & Manuscript & Lean boundary & Location\\
\midrule
\endhead
Newton/$q$-Pochhammer identification
& Paper proof
& Finite $q$-Pochhammer API exists; no report-specific $N_k$ theorem found
& \Cref{gq:chap:q-newton}\\
Iterated $D_q$ and divided differences
& Paper proof
& No one-to-one declaration found
& \Cref{gq:thm:Dq-divided-difference}\\
Gaussian Pascal basis pair
& Paper proof
& Gaussian reciprocity/inversion ingredients formalized; basis pair not packaged
& \Cref{gq:thm:gaussian-pascal-basis}\\
Complete-homogeneous monomial residual
& Paper proof
& Geometric residual moments formalized; arbitrary evaluation remainder not packaged
& \Cref{gq:thm:monomial-residual}\\
Analytic infinite Newton convergence
& Paper proof
& No one-to-one declaration found
& \Cref{gq:thm:analytic-newton-convergence}\\
Endpoint Lebesgue formula and limit
& Paper proof
& Exact finite rational $\ell^1$ identity formalized; real infinite-product limit not
& \Cref{gq:thm:endpoint-lebesgue}\\
Boundary profile and global Lebesgue asymptotic
& Paper proof
& No one-to-one declaration found
& \Cref{gq:thm:lebesgue-profile,gq:thm:global-lebesgue}\\
Newton--Jackson formula
& Paper proof
& No one-to-one declaration found
& \Cref{gq:thm:jackson-newton-moments}\\
Fabius half-base Gaussian transform
& Paper proof
& Inverse-dyadic moment and Gaussian ingredients formalized; transform not packaged
& \Cref{gq:thm:fabius-gaussian-transform}\\
Fabius residue and fixed-jet bounds
& Paper proof
& No one-to-one declarations found
& \Cref{gq:thm:fabius-boundary-residue,gq:thm:fabius-fixed-jets}\\
Fabius gap upper bound
& Paper proof
& No one-to-one declaration found
& \Cref{gq:thm:fabius-gap-upper}\\
Two-sided Fabius gap asymptotic
& Conjecture
& Unproved
& \Cref{gq:conj:fabius-gap-asymptotic}\\
Second-order Lebesgue expansion
& Conjecture
& Unproved
& \Cref{gq:conj:second-order-lebesgue}\\
Hermite saddle displacement
& Conjecture
& Unproved
& \Cref{gq:conj:hermite-saddle}\\
Dyadic fixed-radius $\RvachevUp$ synthesis
& Paper proof
& Directly formalized by
  \path{normalized_sum_Ioo_rvachevDeconvolvedPolynomial_mul_shifted_rvachevUp}
& \Cref{gq:thm:up-polynomial-synthesis}\\
Geometric q-falling Appell formula and synthesis
& Paper proof
& Directly formalized by
  \path{eval_rvachevDeconvolvedPolynomial_qFallingPower}, together with
  \path{normalized_sum_Ioo_rvachevDeconvolvedPolynomial_mul_shifted_rvachevUp}
  at the dyadic mesh
& \Cref{gq:prop:q-Appell-falling}\\
Generic scalar cardinal-decoder synthesis
& Paper proof
& Directly formalized by
  \path{normalized_sum_Ioo_lagrangeRvachevDecoder_mul_shifted_rvachevUp}
& \Cref{gq:thm:gaussian-Appell-decoder}\\
Geometric closed decoder formula
& Paper proof
& Directly formalized by
  \path{geometric_lagrangeRvachevDecoder_eq}; its nonnegative denominator
  exponent is equivalent to the manuscript's negative-power form under
  \(0<q<1\) and \(c>0\)
& \Cref{gq:thm:gaussian-Appell-decoder}\\
Atomic coefficient factorization and full interpolation loop
& Paper proof
& Directly formalized by
  \path{lagrangeRvachevAtomCoefficient_eq_deconvolved_interpolate} and
  \path{sum_Ioo_lagrangeRvachevAtomCoefficient_mul_shifted_rvachevUp}
& \Cref{gq:cor:interpolation-up-factorization}\\
Nodewise biorthogonality and decoder row sum
& Paper proof
& Directly formalized componentwise by
  \path{normalized_sum_Ioo_lagrangeRvachevDecoder_eval_node} and
  \path{sum_lagrangeRvachevDecoder_eq_one}, and in typed form by
  \path{lagrangeRvachevEncoderMatrix_mul_decoderMatrix} and
  \path{sum_lagrangeRvachevDecoderMatrix_row_eq_one}
& \Cref{gq:thm:gaussian-Appell-biorthogonality}\\
Stochastic encoder and Matrix right inverse
& Paper proof
& Directly formalized by \path{lagrangeRvachevEncoderMatrix_nonneg},
  \path{sum_lagrangeRvachevEncoderMatrix_row_eq_one}, and
  \path{lagrangeRvachevEncoderMatrix_mul_decoderMatrix}; the negative-entry
  conclusion is conditional on the positive-overlap hypothesis in
  \path{exists_lagrangeRvachevDecoderMatrix_entry_neg_of_row_overlap}
& \Cref{gq:thm:gaussian-Appell-biorthogonality}\\
Minimum-variation atomic decoder
& Open problem
& Unproved; no optimization or conditioning theorem found
& \Cref{gq:prob:minimum-variation-decoder}\\
Growing-order Hermite anchoring
& Open problem
& Unproved
& \Cref{gq:prob:optimal-geometric-hermite}\\
\bottomrule
\end{longtable}
\normalsize
